\bookchapter{第二十九章：數字健康——同手機建立健康關係}

\chapterintro{你有冇咁樣嘅體驗：本來淨係想掃5分鐘短影片，結果一抬頭已經過咗兩個鐘？或者夜晚瞓喺牀度掃手機，越掃越精神，最後凌晨兩點先瞓着？手機改變咗我哋嘅生活，但亦喺度悄悄吞噬我哋嘅時間同健康。呢一章唔係等你「戒掉手機」，而係幫你重新攞返對手機嘅主動權。}

\booksection{一、先睇嚇你嘅「螢幕時間」——數字等人清醒}

喺改變之前，先搞清楚現狀。手機入面便有一個好多人從未打開過嘅功能，叫「螢幕使用時間」（iPhone）或「數字健康」（安卓），佢能話畀你聽每日到底喺手機上花咗幾多時間，以及呢啲時間花咗喺邊啲App上。

\begin{itemize}
  \item \textbf{iPhone}：設定 > 螢幕使用時間 > 查看所有活動。你會見到每日嘅解鎖次數、攞起手機嘅次數、每個App嘅使用時長。
  \item \textbf{安卓}：設定 > 數字健康與家長控制。華為叫「健康使用手機」，小米叫「螢幕時間管理」。功能類似，都能見到每日使用數據。
\end{itemize}

\textbf{第一次睇呢個數據嘅人，80\%會被嚇到}。我有個朋友睇完之後發現自己每日喺抖音上花咗4個幾鐘，一週就係28小時——相當於一份兼職嘅時間。呢種「數字衝擊」係改變嘅第一步。

\booksection{二、畀App設定「使用限額」——畀自己裝個限速器}

知問題喺邊之後，就可以有針對性咁設定限制。唔需要一刀切將App刪咗，淨係需要畀佢設個「每日額度」：

\begin{enumerate}
  \item \textbf{iPhone}：設定 > 螢幕使用時間 > App限額 > 添加限額。選中你要限制嘅App（比如抖音），設定每日30分鐘。到時間後App圖標會變灰，打開需要額外輸入密碼。
  \item \textbf{安卓}：設定 > 數字健康 > 應用計時器。選中App後設定每日使用上限，到時間後App圖標會變灰。
  \item \textbf{小技巧}：將限額密碼設成一個你唔常用、唔太好記嘅密碼。太容易輸嘅話，你會忍唔住每次都「續命」，限額就形同虛設喇。
\end{enumerate}

\booksection{三、「勿擾模式」嘅正確用法——唔係關掉通知，係關掉焦慮}

好人嘅手機一日響幾十次通知，每響一次注意力就被打斷一次。你以為你淨係「睇一眼通知」，實際上每次打斷後重新集中注意力需要15-20分鐘。勿擾模式唔係等你「失聯」，而係幫你重新掌控注意力。

\begin{itemize}
  \item \textbf{定時開啟}：設定勿擾模式喺夜晚10點到朝早7點自動開啟。咁樣瞓覺時手機唔會因為一條垃圾短訊將你嘈醒。
  \item \textbf{允許重要聯絡穿透}：喺勿擾模式設定入面便，將「允許來電」設為「收藏聯絡人」或「重複來電」（3分鐘內第二次來電會正常響鈴）。咁樣你唔會錯過家人嘅緊急電話，但亦唔會被推銷電話打擾。
  \item \textbf{工作專注模式}：iPhone嘅「專注模式」可以設自訂場景——比如「工作」模式下淨係允許工作相關App通知，「閲讀」模式下所有通知全靜音。
\end{itemize}

\booksection{四、護眼唔靠App靠習慣——20-20-20法則}

長時間睇手機螢幕，眼乾澀、酸脹係必然嘅。護眼模式（藍光過濾）有一定幫助，但真正有效嘅係改變用眼習慣：

\begin{tcbinfo}{20-20-20法則}
每睇螢幕\textbf{20分鐘}，將視線移到\textbf{20英尺（約6米）}以外嘅地方，保持\textbf{20秒}。呢個簡單嘅動作能有效緩解睫狀肌疲勞。

\end{tcbinfo}

\begin{itemize}
  \item \textbf{開啟護眼模式}：設定 > 顯示 > 護眼模式（或夜覽模式），開啟後螢幕色溫變暖，減少藍光。建議設為日落到日出自動開啟。
  \item \textbf{調整螢幕亮度}：環境暗時螢幕亦要暗。好多人喺黑暗中睇手機亮度調太高，呢個係傷眼最嚴重嘅行為。開啟「自動亮度」並確保佢正常工作。
  \item \textbf{保持距離}：手機距離眼睛至少30釐米。好多人鍾意將手機攞得特別近睇，咁樣睫狀肌一直喺緊張收縮狀態，近視加深好快。
\end{itemize}

\booksection{五、「低頭族」嘅頸椎自救指南}

低頭睇手機時，你嘅脖子承受嘅壓力係正常站立時嘅5-6倍。長期低頭會導致頸椎曲度變直、骨質增生、頭暈手麻。以下係幾個簡單嘅改善方法：

\begin{itemize}
  \item \textbf{將手機舉高}：睇手機時盡量將手機舉到與眼睛平齊嘅高度，而唔係低頭睇放喺大髀上嘅手機。雖然手臂會攰，但脖子會多謝你。
  \item \textbf{每30分鐘活動一下脖子}：慢慢向左轉頭到極限，停3秒，回正；再向右轉，停3秒。重複5次。然後做「米字操」——用下巴喺空中寫一個「米」字。
  \item \textbf{靠牆企3分鐘}：後腦勺、肩胛骨、臀部、腳後跟四點貼牆，每日企3分鐘，能有效糾正含胸低頭嘅體態。
\end{itemize}

\booksection{六、睡眠與手機——臥室入面便最該改嘅習慣}

手機螢幕嘅藍光會抑制褪黑素分泌，等你「越睇越精神」。好多人覺得「我掃住掃住就攰喇」，實際上係被藍光欺騙咗大腦，等你放低手機反而唔容易入睡。以下係改善建議：

\begin{itemize}
  \item \textbf{睡前1小時唔睇手機}：呢個係最有效亦最難做到嘅一條。如果做不到，至少開「夜覽/護眼模式」將藍光降到最低。
  \item \textbf{手機放喺伸手夠唔到嘅地方}：將充電器放喺臥室另一頭，而唔係牀頭櫃。咁樣你既唔會忍唔住攞起嚟掃，鬧鐘響咗你亦要起身先至可以閂掉——一舉兩得。
  \item \textbf{開啟「睡眠模式」}：iPhone嘅「睡眠」專注模式可以自動喺睡眠時段鎖定螢幕、簡化鎖屏；安卓各品牌亦有類似功能。
  \item \textbf{用鬧鐘代替手機}：買一個十幾蚊嘅實體鬧鐘，唔用手機當鬧鐘，你會發現少咗「睇一眼時間然後忍唔住掃一下」嘅陷阱。
\end{itemize}

\booksection{七、資訊焦慮症——學識「選擇性失明」}

我哋每日接收嘅資訊量，係古人一世接收嘅量。但人嘅大腦唔係為咗處理咁多資訊而設計嘅，資訊過載會導致焦慮、注意力渙散、決策疲勞。解決方法唔係「獲取更多資訊」，而係「主動屏蔽」：

\begin{itemize}
  \item \textbf{取消關注}：打開你關注嘅所有公眾號、UP主、博主列表，問自己一個問題：「過去一個月我從中獲得過實際價值嗎？」如果冇，取消關注。資訊渠道貴精唔貴多。
  \item \textbf{關閉非必要App推送}：除咗微信、電話、短訊、日曆之外，其他App嘅通知全部關掉。新聞唔需要推送——你主動去睇嘅時候自然會見到。
  \item \textbf{設立「資訊齋戒日」}：每週選一日（比如週日），全天唔睇手機新聞、唔掃短影片。你會發現天唔會塌落嚟，而且週一精神狀態明顯更好。
\end{itemize}

\begin{tcbpractice}{動手：一週數字健康挑戰}
\begin{enumerate}
  \item 今日：打開螢幕使用時間/數字健康，查看你上週嘅使用數據。
  \item 聽日：畀耗時最長嘅娛樂App設定每日限額（建議30-60分鐘）。
  \item 後日：設定勿擾模式定時開啟（22:00-07:00），允許收藏聯絡人來電。
  \item 第四日：睡前1小時唔睇手機，手機充電器放到臥室另一頭。
  \item 第五日：取消關注5個你唔再睇嘅公眾號或博主。
  \item 第六日：嘗試一成日嘅「資訊齋戒」——唔睇新聞唔掃短影片。
  \item 第七日：查看本週嘅螢幕使用時間報告，同上週對比。
\end{enumerate}

\end{tcbpractice}

\begin{tcbsummary}{本章小結}
\begin{itemize}
  \item 先睇螢幕使用時間數據，認清現狀係改變嘅第一步
  \item 畀娛樂App設限額，畀勿擾模式設定時，重新掌控注意力
  \item 20-20-20法則係最簡單有效嘅護眼習慣
  \item 低頭睇手機時脖子承重5-6倍，舉高手機+定期活動脖子
  \item 睡前1小時唔睇手機，充電器放到伸手夠唔到嘅地方
  \item 資訊焦慮嘅解藥係「主動屏蔽」，唔係「獲取更多」
\end{itemize}

\end{tcbsummary}
